\section{Skill System}
\label{sec:skill-system}

\subsection{Overview}

The Skill System defines how Heroes and Monsters act in combat. 
It is built on a unified model for Skills, Impacts, Status effects and Modifiers, shared by both sides:

\begin{itemize}
	\item \textbf{SkillData} -- base definition: cooldown, range, tags, base damage, target rules.
	\item \textbf{Impacts} -- actions executed on cast / hit: deal damage, apply status, trigger other skills.
	\item \textbf{Status} -- buffs, debuffs, damage-over-time and similar timed effects attached to units.
	\item \textbf{ModifierData} -- generic numeric changes applied by gear, passives, buffs, debuffs, runes, crowns, etc.
\end{itemize}

Heroes and Monsters use the same SkillData and Status/Modifier framework. 
The difference lies in where these Skills come from (Hero archetypes, Modules, enemy definition) and which Modifiers apply.

\subsection{Shared Stat Model}

The Skill System relies on a fixed set of five core stats:

\begin{itemize}
	\item \textbf{STR} \colorsquare{STR} -- physical power (melee damage, strength-based scaling).
	\item \textbf{DEX} \colorsquare{DEX} -- agility (attack speed, accuracy, evasion-style interactions).
	\item \textbf{CON} \colorsquare{CON} -- resilience (life, barrier, physical toughness).
	\item \textbf{INT} \colorsquare{INT} -- mental power (spell power, elemental scaling, complex skills).
	\item \textbf{WIL} \colorsquare{WIL} -- willpower (resistance to control, resource sustain, thematic „will“ effects).
\end{itemize}

All Skill and Modifier calculations ultimately work on this stat set and derived combat stats (damage, crit, armor, resistances, etc.).

% RISK: Any later change to this stat set will ripple into all systems (Heroes, Skills, Gear, Crafting, Balance).

\subsection{Skill Model (Heroes and Enemies)}

\subsubsection*{SkillData}

Each Skill is defined by a SkillData object, including:

\begin{itemize}
	\item skill-family (GroundEffect, SinpleHit, Projectile, Nova, etc.)
	\item base damage (physical/elemental or mixed),
	\item cooldown and cast time,
	\item range and target rules (single target, AoE, line, cone, self),
	\item tags (e.g. melee, projectile, DoT, aura, summon, fire, holy, undead),
	\item scaling hooks (which stats modify what parts of the skill).
	\item effect hooks (impacts/status)
\end{itemize}

\subsubsection*{Impacts and Status ("Effects")}

Technically devided in OnCastEffect, OnHitEffect and EndEffects.

\begin{itemize}
	\item \textbf{Impacts} define what happens when a Skill is cast or hits:
	\begin{itemize}
		\item damage instances,
		\item application of Status effects,
		\item triggering of chained skills (e.g. nova on hit, projectiles on kill).
	\end{itemize}
	\item \textbf{Status} objects represent ongoing effects:
	\begin{itemize}
		\item buffs (e.g. increased armor, haste),
		\item debuffs (e.g. slow, vulnerability),
		\item DoTs (e.g. bleeding, burning),
		\item other persistent states (e.g. taunt, mark).
	\end{itemize}
	\item Status effects apply \textbf{ModifierData} to stats or skills while active.
\end{itemize}

\subsubsection*{ModifierData}

ModifierData is the unified language for numeric changes:

\begin{itemize}
	\item flat additions (e.g. +10 armor),
	\item additive \% increases (e.g. +20\% damage),
	\item multiplicative „more“ multipliers (e.g. 30\% more damage),
	\item caps and clamps (e.g. resistances up to 75\%),
	\item occasional replacements or transformations (e.g. convert physical to fire, replace basic attack).
\end{itemize}

All systems (Gear, Skill Orbits, Runes, Crowns, Research) ultimately output ModifierData.

\subsection{Hero Skill Layers}

Hero power and behaviour in combat is defined through several stacked layers:

\begin{enumerate}
	\item Hero archetype and base stats (STR, DEX, CON, INT, WIL).
	\item Gear (from the Crafting System).
	\item \textbf{Crafted Skills (Skill Modules) in Sockets.}
	\item Socket Orbits (passive progression around each socket).
	\item Hero-wide Runes.
	\item Endgame Crowns.
\end{enumerate}

Details about Hero archetypes, roles and progression (HeroDef, HeroInstance, AI profiles) are described in Section~\ref{sec:heroes-system}. 
This section focuses on how Heroes interface with the Skill System via sockets, modules and orbits.


\subsubsection{Skill Modules and Sockets}

Heroes do not „learn“ traditional skill trees; instead, they use \textbf{Skill Modules} socketed into fixed positions.

\begin{itemize}
	\item Each Hero has \textbf{5 chakra-style sockets}, not bound to gear:
	\begin{itemize}
		\item e.g. Heart, Hands, Head, Core, Movement.
	\end{itemize}
	\item Sockets are \textbf{semantically distinct}, not statistically tied to gear slots:
	\begin{itemize}
		\item e.g. a socket may be associated with crit/burst, sustain, mobility, defence.
	\end{itemize}
	\item Each Scoket has a fix resonance to 2 Stats - only Skills of that stat can be inserted into the Socket\\
	
	\begin{tabular}{llll}
		\textbf{Slot}  & \textbf{allowed Stats incl. Color (Hex)} & \textbf{Theme} \\
		Mind     & \textbf{WIL} \colorsquare{WIL} \texttt{\#26e3e3} + \textbf{INT} \colorsquare{INT} \texttt{\#5626e3} & focus, mental power \\
		Heart       & \textbf{STR} \colorsquare{STR} \texttt{\#e3263e} + \textbf{WIL} \colorsquare{WIL} \texttt{\#26e3e3} & courage, passion, determination \\
		Core       & \textbf{INT} \colorsquare{INT} \texttt{\#5626e3} + \textbf{STR} \colorsquare{STR} \texttt{\#e3263e} & technique, core strength, effective power \\
		Handling   & \textbf{STR} \colorsquare{STR} \texttt{\#e3263e} + \textbf{DEX} \colorsquare{DEX} \texttt{\#9ce326} & Weapon handling, dexterity, precision  \\
		Movement   & \textbf{DEX} \colorsquare{DEX} \texttt{\#9ce326} + \textbf{WIL} \colorsquare{WIL} \texttt{\#26e3e3} & Speed, endurance, reaction \\
		
	\end{tabular}
	\item Each socket can hold \textbf{exactly one Skill Module item} from the Inventory.
	\item Socket progression:
	\begin{itemize}
		\item start with 1 active socket at low level,
		\item additional sockets unlock as the Hero levels up (e.g. at level milestones).
	\end{itemize}
	\item When a Skill Module is socketed:
	\begin{itemize}
		\item the Hero gains a SkillData instance that is \textbf{fully defined by the module itself}:
		\begin{itemize}
			\item \textbf{Skill identity:} reference to a concrete Skill (SkillData / SkillDef),
			\item \textbf{Skill Tier/Level:} controls base numbers / scaling budget,
			\item \textbf{optional DamageType conversion:} applied as a defined modifier (e.g. phys $\rightarrow$ fire).
		\end{itemize}
		\item The Hero archetype does \textbf{not} transform the skill into a different one.
		Instead, archetypes influence performance through passives, tags, stats, AI behaviour and global modifiers.
	\end{itemize}
\end{itemize}

\noindent
\textbf{Note:} How Skill Modules are crafted (selecting Skill + Tier/Level + DamageType conversion) is defined in the Crafting chapter.
This section treats Skill Modules as existing inventory items.

\paragraph{Module Role}

Skill Modules represent the main active skill layer for Heroes:

\begin{itemize}
	\item they define \textbf{what} the skill does (projectile, aura, buff, summon, etc.) by referencing a concrete SkillData,
	\item they define \textbf{how strong} the base skill is via Tier/Level,
	\item they optionally define a \textbf{DamageType conversion} (if present),
	\item Orbits, Gear, Runes and Crowns further modify behaviour and scaling via tags and ModifierData.
\end{itemize}

\begin{minipage}[t]{0.48\textwidth}
	\textbf{Element Prefixes}\\
	
	\begin{tabular}{r| l}
		\textbf{DMG-Type} & \textbf{Prefix} \\
		\hline
		Fire      & \glqq \textit{Blazing}\grqq   \\
		Cold      & \glqq \textit{Glacial}\grqq   \\
		Lightning & \glqq \textit{Static}\grqq    \\
		Poison    & \glqq \textit{Venomous}\grqq  \\
		Demonic   & \glqq \textit{Infernal}\grqq  \\
		Holy      & \glqq \textit{Radiant}\grqq   \\
	\end{tabular}
\end{minipage}
\hfill

\vspace{1em}


% RISK: If Modules are too numerous or too similar, the system loses clarity.
% TODO: For MVP, keep Module count per archetype low and strongly themed.

\subsubsection{Socket Orbits (Passive Progression)}

Each socket has its own \textbf{Orbit}: a layered progression ring around the socket that works similar to electron shells in an atom model:

\begin{itemize}
	\item the \emph{socket} is the core,
	\item each \emph{orbit tier} is an electron shell around that core,
	\item orbit \emph{points} are the electrons on that shell.
	\item an orbit \emph{Mod} is a "Super-electron" (not scientific related) unlocked when electron hsell full
\end{itemize}

\paragraph{Structure}

\begin{itemize}
	\item Each socket orbit contains a fixed number of points (e.g. 22 points) spread across multiple tiers (Tier~1, Tier~2, Tier~3).
	\item Points are \textbf{small, incremental stat buffs}:
	\begin{itemize}
		\item e.g. +\% damage, +\% life, +crit chance, +status effect, +resource regen,
		\item tuned per socket identity (crit/burst vs. sustain vs. utility, etc.).
	\end{itemize}
	\item Points are invested in free order along the orbit path for that socket.
\end{itemize}

\paragraph{Tier Gates and Orbit Mods}

Orbit progression is divided into tiers, each with a gate:

\begin{itemize}
	\item After investing a required number of points on the current tier (e.g. 6--8 points on Tier~1),
	the player unlocks a \textbf{Tier Mod Choice} for that tier.
	\item A Tier Mod is a \textbf{behaviour/synergy modifier}, not a small stat buff. Examples:
	\begin{itemize}
		\item convert a single-target skill into a small AoE,
		\item add a nova or chain effect on hit,
		\item add a DoT component or a mark on hit,
		\item convert damage type (e.g. phys to fire) or add penetration,
		\item change targeting logic (e.g. prioritise bosses / clusters).
	\end{itemize}
	\item When a Tier Mod is chosen:
	\begin{itemize}
		\item its effect is applied permanently to the socketed skill,
		\item the next outer orbit tier is unlocked and becomes available for point investment.
	\end{itemize}
\end{itemize}

\paragraph{Socket Identity}

Each socket has its own themed orbit and Tier Mod pools:

\begin{itemize}
	\item \textbf{Example (conceptual only):}
	\begin{itemize}
		\item Torso socket: sustain/defence (life, barrier, mitigation, „on damage taken“ effects).
		\item Right Hand: crit and burst (crit chance/multi, „execute“ style effects, high-risk/high-reward).
		\item Left Hand: utility/control (crowd control, debuffs, resource manipulation).
		\item Head: AoE/range/complex mechanics (larger areas, extra projectiles, complex triggers).
	\end{itemize}
	\item Tier Mods are drawn from socket-specific pools to reinforce the identity of that slot.
\end{itemize}

\paragraph{Progression Source}

Orbit points are gained through levle-ups of the Heroe:

\begin{itemize}
	\item Hero level ups,
	\item Research unlocks (Skill Tiers in the Heroes Tree),
	\item possibly specific achievements or milestones (future).
\end{itemize}

The intended pattern is:

\begin{itemize}
	\item frequent \textbf{small numeric gains} through orbit points,
	\item occasional \textbf{big behavioural steps} when a Tier Mod is unlocked,
	\item each unlocked Tier Mod both upgrades the current skill and opens access to the next orbit tier.
\end{itemize}

% RISK: Too many small points per orbit can feel grindy; tier mod unlock thresholds must be tuned to feel meaningful.
% TODO: Define exact points-per-tier and thresholds per socket in the Heroes chapter / balance sheets.

\subsubsection{Crowns (Endgame)}

Crowns are an \textbf{endgame progression layer}:

\begin{itemize}
	\item Crown sockets are unlocked via the Heroes Research Tree (see Research System).
	\item Crowns apply \textbf{large, rare, endgame-level modifiers}, such as:
	\begin{itemize}
		\item transforming a core skill in a fundamental way,
		\item adding strong, unique effects with clear downsides,
		\item pushing certain stats beyond normal caps.
	\end{itemize}
	\item Crowns are intentionally rare and limited; they are not part of the early game.
\end{itemize}

For MVP, Crowns can be treated as \emph{future or late} content and do not need full implementation.

% RISK: Too many overlapping layers (Orbits, Runes, Crowns, Gear) can overwhelm players.
% TODO: Crowns should be reserved for clearly signposted endgame goals.

\subsection{Skill Loadout and Combat Usage}

\subsubsection*{Hero Loadout}

A Hero's combat loadout is defined by:

\begin{itemize}
	\item the number of \textbf{active sockets} (unlocked by level),
	\item which \textbf{Modules} are slotted into those sockets,
	\item which \textbf{Orbit mods} have been chosen around each socket,
	\item hero-wide \textbf{Runes} and possible \textbf{Crowns},
	\item equipped Gear and its interactions with skill tags.
\end{itemize}

In practice:

\begin{itemize}
	\item early game: 1 active socket → 1 major skill per Hero,
	\item mid game: 2--3 sockets → 2--3 major skills per Hero,
	\item late game: 3--4 (5) sockets with fully developed Orbits.
\end{itemize}

\subsubsection*{AI-Driven Skill Usage}

There is no direct player control over skill priority in combat:

\begin{itemize}
	\item each Hero follows an \textbf{archetype-specific AI profile}:
	\begin{itemize}
		\item tanks prioritise taunts/defensive skills and high-HP targets,
		\item damage dealers prioritise high-value or low-HP targets,
		\item support archetypes prioritise buffs and heals.
	\end{itemize}
	\item AI rules consider:
	\begin{itemize}
		\item skill tags (e.g. „heal“, „buff“, „AoE nuke“),
		\item cooldown states,
		\item simple conditions (e.g. ally HP, enemy clustering).
	\end{itemize}
\end{itemize}

The player expresses build intent primarily through:

\begin{itemize}
	\item \textbf{skill choices} (which concrete skills are socketed),
	\item orbit mods (how those skills behave),
	\item gear and runes (which stats and tags are favoured),
	\item optional damage type conversion on the skill module.
\end{itemize}

% RISK: Lack of explicit skill priority control can frustrate highly engaged players.
% TODO: Evaluate later whether a simple „priority list“ per Hero is needed, once AI is stable.

\subsection{Interactions with Other Systems}

\subsubsection*{Research System}

\begin{itemize}
	\item The \textbf{Heroes Research Tree} unlocks:
	\begin{itemize}
		\item new Hero archetypes (Hero packs),
		\item deeper Skill tiers (access to higher Orbit tiers),
		\item Crown sockets as endgame goals.
	\end{itemize}
	\item Research does not grant direct one-run buffs, but expands long-term build options.
\end{itemize}

\subsubsection*{Gear and Crafting}

\begin{itemize}
	\item Gear affects Skills primarily via:
	\begin{itemize}
		\item stats that scale skills (STR, INT, crit, penetration),
		\item tag-based bonuses (e.g. increased projectile skill power, increased AoE, reduced cooldown).
	\end{itemize}
	\item The Crafting System defines how gear is created and upgraded; Skill System uses the resulting stats and tags.
\end{itemize}

\subsubsection*{Inventory and Modules}

\begin{itemize}
	\item Skill Modules and Runes are inventory items:
	\begin{itemize}
		\item Modules and Runes are managed in the Hideout inventory screen,
		\item equipping is done via sockets on the Hero panel.
	\end{itemize}
	\item No inventory management is available during Maps; skill loadouts are fixed for the duration of a run.
\end{itemize}

\subsubsection*{Map \& Combat System}

\begin{itemize}
	\item The Map \& Combat System calls into the Skill System to:
	\begin{itemize}
		\item execute Skills for Heroes and Monsters,
		\item resolve Impacts and Status effects,
		\item track kills, damage and other metrics for Research.
	\end{itemize}
\end{itemize}

\subsection{MVP Scope for the Skill System}

For the MVP, the Skill System should deliver:

\begin{itemize}
	\item A stable \textbf{SkillData / Status / Modifier} framework used by Heroes and Monsters.
	\item A limited set of \textbf{Hero archetypes} (e.g. 3--5) with:
	\begin{itemize}
		\item clear tag-based identities (what tags they amplify),
		\item a small starter pool of \textbf{crafted skills} that fit their theme,
		\item basic Orbit layouts (Tier I/II, Tier III optional).
	\end{itemize}
	\item \textbf{Sockets and Modules}:
	\begin{itemize}
		\item start with 1 active socket per Hero,
		\item unlock up to 2--3 sockets for MVP,
		\item small, strongly themed \textbf{skill pool} (skills are concrete SkillData items).
	\end{itemize}
	\item \textbf{Orbits}:
	\begin{itemize}
		\item functional Tier I/II for each socket,
		\item Tier III and complex mutations may be endgame or future additions.
	\end{itemize}
	\item \textbf{Runes}:
	\begin{itemize}
		\item optional, simple hero-wide bonuses (1 slot is sufficient for MVP),
		\item complex interactions can be postponed.
	\end{itemize}
	\item \textbf{Crowns}:
	\begin{itemize}
		\item treated as late-game or future content,
		\item only minimal scaffolding needed in MVP.
	\end{itemize}
\end{itemize}

% TODO: Finalise MVP archetype list and per-archetype Module/Orbit designs in a dedicated Heroes chapter.
